\section{Reinforcement Learning with GRPO}
\label{sec:grpo}

\subsection{Motivation}
\label{sec:grpo:motivation}

While SFT teaches the model domain-specific formatting and reasoning patterns, it is fundamentally limited to imitating the training data. On questions the base model already solves correctly, SFT reinforces existing capabilities; on questions where the model fails, SFT provides the correct answer but does not teach the model \emph{how to arrive at it} from its own reasoning process. This limitation is reflected in our results: the SFT model trained on the 600 questions where GRPO cannot learn (those the base model solves either 0/4 or 4/4 times, producing zero reward variance) shows no improvement and even slight regression on the remaining 266 mixed-difficulty questions (Section~\ref{sec:grpo:results}).

Group Relative Policy Optimization (GRPO)~\cite{shao2024deepseekmath} offers a complementary approach: rather than imitating expert solutions, the model generates multiple candidate answers and learns to prefer reasoning chains that lead to correct results. For reliability engineering, where answers are verifiable numerical values, this provides a natural reward signal without requiring a learned reward model.

\subsection{GRPO Algorithm}
\label{sec:grpo:method}

GRPO generates $G$ completions $\{o_1, \ldots, o_G\}$ for each prompt $q$, scores each with a reward function $r(o_i, q)$, and computes \emph{group-relative} advantages:
\begin{equation}
    \hat{A}_i = \frac{r(o_i, q) - \bar{r}}{\sigma_r + \epsilon},
    \label{eq:grpo_advantage}
\end{equation}
where $\bar{r}$ and $\sigma_r$ are the mean and standard deviation of rewards within the group. The policy is updated to maximize:
\begin{equation}
    \mathcal{L}_{\text{GRPO}} = -\frac{1}{G} \sum_{i=1}^{G} \hat{A}_i \cdot \frac{\log \pi_\theta(o_i \mid q)}{|o_i|},
    \label{eq:grpo_loss}
\end{equation}
where $|o_i|$ is the response length in tokens, normalizing the loss to prevent bias toward shorter completions (Dr.~GRPO~\cite{deepseek2025r1}).

Unlike PPO, GRPO requires no critic model: the group mean serves as the baseline. Unlike DPO, it does not require pre-computed preference pairs but generates them online. This makes GRPO particularly suited to domains with verifiable rewards and limited data.

\subsection{Dynamic Sampling (DAPO)}
\label{sec:grpo:dapo}

A key challenge with GRPO on small datasets is that many prompts produce zero reward variance: either all $G$ generations are correct (no learning signal for improvement) or all are wrong (no positive signal). Following DAPO (Dynamic sampling Policy Optimization)~\cite{yu2025dapo}, we implement \emph{dynamic sampling}: when all $G$ generations produce the same reward, the prompt is discarded and a new one is sampled. This ensures every training step provides a meaningful gradient signal.

We track \emph{useful steps} (prompts with reward variance) separately from \emph{wasted steps} (prompts discarded due to zero variance). In practice, dynamic sampling proves essential: to accumulate 200 useful steps, the model required 260 total attempts (60 wasted, 23.1\% waste rate), meaning nearly \textbf{one in four prompts} produced zero-variance rewards and would have contributed no gradient without resampling. Without dynamic sampling, these wasted steps would either inject pure noise (all-correct prompts reinforce the status quo) or provide no learning signal at all (all-wrong prompts), substantially degrading training efficiency on a 266-question dataset.

\subsection{Reward Design}
\label{sec:grpo:reward}

Our reward function uses the same automated numerical comparison as the SFT evaluation pipeline (Section~\ref{sec:sft:eval}). Given a model response $o$ and ground-truth target $y$, we extract the last \verb|\boxed{}| value $\hat{y}$ and compute a graded reward based on relative error $\epsilon = |\hat{y} - y| / (|y| + 10^{-9})$:

\begin{equation}
    r(o, q) = \begin{cases}
        1.0 & \text{if } \epsilon \leq 0.1\% \\
        0.8 & \text{if } \epsilon \leq 1\% \\
        0.4 & \text{if } \epsilon \leq 5\% \\
        0.01 & \text{if } \epsilon > 5\% \\
        0.0 & \text{if no } \texttt{\textbackslash boxed\{\}} \text{ found} \\
        -0.3 & \text{if } > 2 \text{ boxed values}
    \end{cases}
    \label{eq:reward}
\end{equation}

The graded structure provides a smoother reward landscape than binary correct/incorrect: partially correct answers ($\epsilon \leq 5\%$) receive positive reward, encouraging the model to refine its computation rather than abandon an approach entirely. The penalty for multiple \verb|\boxed{}| values discourages hedging.

\subsection{Training Data Split}
\label{sec:grpo:split}

The 866-question dataset (v4) is split based on base-model screening: each question is evaluated 4 times with the base Qwen3-8B model, and questions are categorized by success rate:

\begin{itemize}
    \item \textbf{600 questions (0/4 or 4/4 correct):} Used for SFT training. These include 210 questions the base model always fails (0/4) and 390 it always solves (4/4). Both extremes produce zero reward variance under GRPO---either all generations are correct or all are wrong---so they cannot provide a gradient signal for RL. However, they are suitable for supervised learning from ground-truth answers.
    \item \textbf{266 questions (1/4 to 3/4 correct):} Used for GRPO training. These have natural reward variance---some generations succeed, others fail---making them ideal for group-relative optimization.
\end{itemize}

This split ensures zero data leakage between SFT and GRPO stages and targets GRPO training at the model's learning frontier---questions where the base model sometimes succeeds and sometimes fails, providing the contrastive signal that GRPO requires.

\subsection{Training Configuration}
\label{sec:grpo:config}

We use a custom training loop (no TRL dependency) with $G\!=\!4$ generations per prompt, 200 useful steps, LR=$1\!\times\!10^{-5}$ with cosine decay, gradient accumulation of 4, and $\beta\!=\!0$ (no KL penalty). The pipeline loads Qwen3-8B, applies the SFT fold\_4 LoRA ($r\!=\!16$, $\alpha\!=\!32$, $\sim$43.6M trainable parameters), and runs GRPO with dynamic sampling on the 266 mixed questions. Full configuration details are in Appendix~A.

\subsection{GRPO Results}
\label{sec:grpo:results}

\subsubsection{In-Distribution Evaluation (266 Mixed Questions)}

Table~\ref{tab:grpo_266q} presents accuracy on the 266 GRPO training questions across all model configurations. All evaluations use deterministic seeding (seed $= 3407 + q_i$) with temperature 1.0 and top-$p$ 0.95, matching the GRPO generation settings.

\begin{table}[t]
\centering
\small
\caption{Accuracy on 266 mixed questions. $\Delta$ and $p$ (McNemar) are vs.\ SFT fold\_4.}
\label{tab:grpo_266q}
\begin{tabular}{@{}lcccc@{}}
\toprule
Model & Acc. & Trunc. & $\Delta$ & $p$ \\
\midrule
Qwen3-8B base              & 50.4\% & 4  & --- & --- \\
SFT fold\_4                & 50.8\% & 8  & ref. & --- \\
GRPO from scratch, step 80 & 50.8\% & 1 & $+$0.0 & 1.00 \\
\midrule
GRPO from SFT, step 200    & 52.3\% & 6  & $+$1.5 & --- \\
GRPO from SFT, step 100    & 54.1\% & 9  & $+$3.3 & --- \\
\textbf{GRPO from SFT, step 80} & \textbf{57.9\%} & 10 & $\boldsymbol{+7.1}$ & 0.056 \\
\bottomrule
\end{tabular}
\end{table}

The best configuration (GRPO from SFT at step~80) achieves \textbf{57.9\%} accuracy, $+7.1$pp over the SFT fold\_4 starting point (McNemar $p=0.056$, approaching significance with 54 wins vs.\ 35 losses). Performance degrades with continued training: $-3.8$pp from step~80 to step~100 and $-5.6$pp to step~200, indicating overfitting after approximately 80 useful steps on this 266-question dataset.

\subsubsection{Out-of-Distribution and Ablation Results}

Table~\ref{tab:grpo_holdout_ablation} presents two key analyses: (a)~generalization on 54 independent holdout questions never seen during SFT or GRPO training, and (b)~the effect of SFT warm-start on GRPO effectiveness.

\begin{table}[t!]
\centering
\small
\caption{(a) Holdout accuracy (54 unseen questions) and (b) SFT warm-start ablation (step 80, 266q). $p$: McNemar test.}
\label{tab:grpo_holdout_ablation}
\begin{tabular}{@{}lcccc@{}}
\toprule
\multicolumn{5}{@{}l}{\textit{(a) Independent holdout (54 questions)}} \\
\midrule
Model & Cor. & Acc. & $\Delta$ & $p$ \\
\midrule
Qwen3-8B base & 29 & 53.7\% & --- & ---    \\
SFT fold\_4   & 18 & 33.3\% & $-$20.4 & 0.019 \\
GRPO from SFT, step 80 & 29 & 53.7\% & $+$0.0 & 0.804 \\
GRPO from scratch, step 200 & 30 & 55.6\% & $+$1.9 & 1.00 \\
\midrule
\multicolumn{5}{@{}l}{\textit{(b) SFT warm-start ablation (266q, step 80)}} \\
\midrule
Init & Acc. & $\Delta$ & Waste & $p$ \\
\midrule
Fresh LoRA & 50.8\% & ref. & 29.3\% & --- \\
\textbf{SFT fold\_4} & \textbf{57.9\%} & $\boldsymbol{+7.1}$ & 23.1\% & 0.104 \\
\bottomrule
\end{tabular}
\end{table}

\paragraph{Holdout generalization.} GRPO from SFT at step~80 achieves 53.7\%, matching the base model exactly ($p=0.804$). While the in-distribution gain does not transfer, crucially it does not regress either. The SFT fold\_4 alone drops to 33.3\% ($-20.4$pp, $p=0.019$), a \emph{statistically significant} degradation. GRPO fully recovers this catastrophic forgetting ($p=0.023$ vs.\ SFT, 11 wins vs.\ 2 losses). GRPO from scratch at step~200 achieves 55.6\% ($+1.9$pp), though this amounts to a single additional question.

\paragraph{Caveat on holdout size.} With only 54 questions, each question represents $\sim$1.85pp of accuracy, limiting statistical power. We deliberately chose not to augment the holdout via paraphrasing: while this would increase sample size, it would introduce repeated mathematical content and risk inflating accuracy through surface-level memorization. A larger holdout composed of \emph{distinct} problems would strengthen these findings.

\paragraph{SFT warm-start ablation.} The SFT warm-start is essential: GRPO from scratch barely exceeds the base model ($+0.4$pp, likely within noise), while GRPO from the SFT checkpoint gains $+7.5$pp. The SFT initialization also produces a lower waste rate (23.1\% vs.\ 29.3\%), indicating more diverse reward signals and more efficient dynamic sampling. This is paradoxical: the SFT alone degrades holdout performance by $-20.4$pp, yet it provides a critical initialization for GRPO. We hypothesize that SFT teaches domain-specific \emph{format and vocabulary} (reliability formulas, structured computation) without teaching correct \emph{application} to hard problems. GRPO then leverages this structural foundation to optimize for numerical correctness. Without SFT, the LoRA starts at zero and must simultaneously learn both format and reasoning---too many degrees of freedom for 266 questions.

We also trained with LR=$5\!\times\!10^{-6}$ and $2\!\times\!10^{-5}$; training dynamics (waste rates 20.9--23.1\%) were similar across all three learning rates. Full results are in Appendix~A.
